% problem_id: S001-lemma-finite-surjective-etale-cover
% source_id: S001 (Stacks Project)
% source_file: spaces-more-groupoids.tex
% statement_locator: spaces-more-groupoids.tex lines 1630--1637
% proof_locator: spaces-more-groupoids.tex lines 1639--1792
% env: lemma
% id_note: label 在本来源内唯一
% pairing: tight (verified)
% status: complete (statement + author proof verbatim + reason + locators)
%
\begin{lemma}
\label{lemma-finite-surjective-etale-cover}
Let $S$ be a scheme.
Let $f : X \to Y$ be a morphism of algebraic spaces over $S$.
If $f$ is separated and locally quasi-finite, then there exists a
scheme $U$ \'etale over $Y$ and a surjective \'etale morphism
$U \to (X/Y)_{fin}$ over $Y$.
\end{lemma}

\begin{proof}
Note that the assertion makes sense by the result of
Lemma \ref{lemma-finite-diagonal}
on the diagonal of $(X/Y)_{fin}$, see
Spaces, Lemma \ref{spaces-lemma-representable-diagonal}.
Let $V$ be a scheme and let $V \to Y$ be a surjective \'etale morphism. By
Lemma \ref{lemma-finite-pullback}
the morphism $(V \times_Y X/V)_{fin} \to (X/Y)_{fin}$ is
a base change of the map $V \to Y$ and hence is surjective and \'etale, see
Spaces,
Lemma \ref{spaces-lemma-base-change-representable-transformations-property}.
Hence it suffices to prove the lemma for $(V \times_Y X/V)_{fin}$.
(Here we implicitly use that the composition of representable, surjective, and
\'etale transformations of functors is again representable, surjective, and
\'etale, see
Spaces, Lemmas \ref{spaces-lemma-composition-representable-transformations} and
\ref{spaces-lemma-composition-representable-transformations-property}, and
Morphisms, Lemmas \ref{morphisms-lemma-composition-surjective} and
\ref{morphisms-lemma-composition-etale}.)
Note that the properties of being separated and locally quasi-finite
are preserved under base change, see
Morphisms of Spaces,
Lemmas \ref{spaces-morphisms-lemma-base-change-separated} and
\ref{spaces-morphisms-lemma-base-change-quasi-finite}.
Hence $V \times_Y X \to V$ is separated and locally quasi-finite as well,
and by
Morphisms of Spaces, Proposition
\ref{spaces-morphisms-proposition-locally-quasi-finite-separated-over-scheme}
we see that $V \times_Y X$ is a scheme as well.
Thus we may assume that $f : X \to Y$ is a separated and locally quasi-finite
morphism of schemes.

\medskip\noindent
Pick a point $y \in Y$. Pick $x_1, \ldots, x_n \in X$ points
lying over $y$. Pick an \'etale neighbourhood $a : (U, u) \to (Y, y)$ and a
decomposition
$$
U \times_S X =
W \amalg
\ \coprod\nolimits_{i = 1, \ldots, n}
\ \coprod\nolimits_{j = 1, \ldots, m_j}
V_{i, j}
$$
as in
More on Morphisms, Lemma
\ref{more-morphisms-lemma-etale-splits-off-quasi-finite-part-technical-variant}.
Pick any subset
$$
I \subset \{(i, j) \mid 1 \leq i \leq n, \ 1 \leq j \leq m_i\}.
$$
Given these choices we obtain a pair $(a, Z)$ with
$Z = \bigcup_{(i, j) \in I} V_{i, j}$
which satisfies conditions \ref{equation-finite-conditions}. In other words
we obtain a morphism $U \to (X/Y)_{fin}$. The construction of this morphism
depends on all the things we picked above, so we should really write
$$
U(y, n, x_1, \ldots, x_n, a, I) \longrightarrow (X/Y)_{fin}
$$
This morphism is \'etale by Lemma \ref{lemma-finite-criterion-etale}.

\medskip\noindent
Claim: The disjoint union of all of these is surjective onto $(X/Y)_{fin}$.
It is clear that if the claim holds, then the lemma is true.

\medskip\noindent
To show surjectivity we have to show the following (see
Spaces, Remark \ref{spaces-remark-warning}): Given a scheme $T$ over
$S$, a point $t \in T$, and a map $T \to (X/Y)_{fin}$ we can find a datum
$(y, n, x_1, \ldots, x_n, a, I)$ as above such that
$t$ is in the image of the projection map
$$
U(y, n, x_1, \ldots, x_n, a, I) \times_{(X/Y)_{fin}} T \longrightarrow T.
$$
To prove this we may clearly replace $T$ by
$\Spec(\overline{\kappa(t)})$ and
$T \to (X/Y)_{fin}$ by the composition
$\Spec(\overline{\kappa(t)}) \to T \to (X/Y)_{fin}$.
In other words, we may assume that $T$ is
the spectrum of an algebraically closed field.

\medskip\noindent
Let $T = \Spec(k)$ be the spectrum of an algebraically closed
field $k$. The morphism $T \to (X/Y)_{fin}$ is given by a pair
$(T \to Y, Z)$ satisfying conditions \ref{equation-finite-conditions}.
Here is a picture:
$$
\xymatrix{
& Z \ar[d] \ar[r] & X \ar[d] \\
\Spec(k) \ar@{=}[r] & T \ar[r] & Y
}
$$
Let $y \in Y$ be the image point of $T \to Y$.
Since $Z$ is finite over $k$ it has finitely many points.
Thus there exist finitely many points $x_1, \ldots, x_n \in X$
such that the image of $Z$ in $X$ is contained in $\{x_1, \ldots, x_n\}$.
Choose $a : (U, u) \to (Y, y)$ adapted to $y$ and $x_1, \ldots, x_n$ as
above, which gives the diagram
$$
\xymatrix{
W \amalg
\ \coprod\nolimits_{i = 1, \ldots, n}
\ \coprod\nolimits_{j = 1, \ldots, m_j}
V_{i, j} \ar[d] \ar[r] &
X \ar[d] \\
U \ar[r] & Y.
}
$$
Since $k$ is algebraically closed and
$\kappa(y) \subset \kappa(u)$ is finite separable
we may factor the morphism
$T = \Spec(k) \to Y$ through the morphism
$u = \Spec(\kappa(u)) \to \Spec(\kappa(y)) = y \subset Y$.
With this choice we obtain the commutative diagram:
$$
\xymatrix{
Z \ar[d] \ar[r] &
W \amalg
\ \coprod\nolimits_{i = 1, \ldots, n}
\ \coprod\nolimits_{j = 1, \ldots, m_j}
V_{i, j} \ar[d] \ar[r] &
X \ar[d] \\
\Spec(k) \ar[r] &
U \ar[r] & Y
}
$$
We know that the image of the left upper arrow ends up in
$\coprod V_{i, j}$. Recall also that $Z$ is an open subscheme
of $\Spec(k) \times_Y X$ by definition of $(X/Y)_{fin}$
and that the right hand square is a fibre product square.
Thus we see that
$$
Z \subset
\coprod\nolimits_{i = 1, \ldots, n}\ \coprod\nolimits_{j = 1, \ldots, m_j}
\Spec(k) \times_U V_{i, j}
$$
is an open subscheme. By construction (see
More on Morphisms, Lemma
\ref{more-morphisms-lemma-etale-splits-off-quasi-finite-part-technical-variant})
each $V_{i, j}$ has a unique point $v_{i, j}$ lying over $u$
with purely inseparable residue field extension
$\kappa(v_{i, j})/\kappa(u)$. Hence each
scheme $\Spec(k) \times_U V_{i, j}$ has exactly one
point. Thus we see that
$$
Z = \coprod\nolimits_{(i, j) \in I} \Spec(k) \times_U V_{i, j}
$$
for a unique subset
$I \subset \{(i, j) \mid 1 \leq i \leq n, \ 1 \leq j \leq m_i\}$.
Unwinding the definitions this shows that
$$
U(y, n, x_1, \ldots, x_n, a, I) \times_{(X/Y)_{fin}} T
$$
with $I$ as found above is nonempty as desired.
\end{proof}
